% Generated draft: requires mathematical and source review.
\documentclass[11pt]{article}
\usepackage[margin=25mm]{geometry}
\usepackage{fontspec}
\setmainfont{Latin Modern Roman}
\usepackage{amsmath,amssymb,mathtools,unicode-math}
\setmathfont{Latin Modern Math}
\usepackage{enumitem,needspace,xurl,hyperref}
\setlength{\parindent}{0pt}
\setlength{\parskip}{5pt}
\setlength{\emergencystretch}{3em}
\newcommand{\N}{\mathbb{N}}
\newcommand{\Z}{\mathbb{Z}}
\newcommand{\Q}{\mathbb{Q}}
\newcommand{\R}{\mathbb{R}}
\newcommand{\C}{\mathbb{C}}
\newcommand{\F}{\mathbb{F}}
\newcommand{\A}{\mathbb{A}}
\newcommand{\PP}{\mathbb{P}}
\newcommand{\E}{\mathbb{E}}
\newcommand{\eps}{\varepsilon}
\newcommand{\dd}{\,\mathrm{d}}
\begin{document}
\section*{UnsolvedMath ID: 30003023}
\textbf{Source problem code:} catalogue entry
\subsection{Definitions and mathematical statement}
Use the standard pure field language \(L_{\mathrm{ring}}=\{0,1,+,-,\cdot\}\). A complete first-order theory \(T\) has the independence property if there is a formula \(\varphi(\bar x;\bar y)\) such that, for every integer \(n\geq 1\), some model \(M\models T\) contains tuples \(\bar a_1,\ldots,\bar a_n\) and \(\bar b_S\), one for each \(S\subseteq\{1,\ldots,n\}\), satisfying
\[
M\models\varphi(\bar a_i;\bar b_S)\quad\Longleftrightarrow\quad i\in S.
\]
The theory is \emph{NIP} if it does not have the independence property. A structure is NIP when its complete theory is NIP.

A valuation on a field \(K\) is a map \(v:K\to\Gamma\cup\{\infty\}\), where \(\Gamma\) is an ordered abelian group, such that \(v(0)=\infty\), the restriction \(K^\times\to\Gamma\) is surjective, and
\[
v(ab)=v(a)+v(b),\qquad v(a+b)\geq\min\{v(a),v(b)\}.
\]
Here \(\infty\) is larger than every element of \(\Gamma\), with \(\infty+\gamma=\infty\). Write
\[
\mathcal O_v=\{a\in K:v(a)\geq0\},\qquad
\mathfrak m_v=\{a\in K:v(a)>0\},\qquad
Kv=\mathcal O_v/\mathfrak m_v.
\]
The valuation is \emph{henselian} if, for every monic polynomial \(f\in\mathcal O_v[X]\) and every \(\alpha\in Kv\) with \(\overline f(\alpha)=0\) and \(\overline f'(\alpha)\neq0\), there is \(a\in\mathcal O_v\) reducing to \(\alpha\) such that \(f(a)=0\). Interpret the valued field as the structure \((K,\mathcal O_v)\) in \(L_{\mathrm{val}}=L_{\mathrm{ring}}\cup\{\mathcal O\}\), with a unary predicate for the valuation ring.

\textbf{Question.} For every field \(K\) and every henselian valuation \(v\) on \(K\), does
\[
\operatorname{Th}_{L_{\mathrm{ring}}}(K)\text{ is NIP}
\quad\Longrightarrow\quad
\operatorname{Th}_{L_{\mathrm{val}}}(K,\mathcal O_v)\text{ is NIP}
\]
hold? The supplied statement imposes no characteristic restriction and does not exclude the trivial valuation.

\textbf{Unresolved qualifications.} The source record does not explicitly specify the languages; the formulation above uses the standard pure field language and valuation-ring predicate convention. Its supplied literature assessment instead discusses henselian expansions with NIP residue field in unspecified remaining imperfect or separably closed cases. That wording does not state the same hypothesis as NIP of \(K\), and the record supplies no equivalence or precise description of those remaining cases. It is retained as a qualification of the supplied literature assessment, not as a replacement or restriction of the question.
\subsection{Short English statement}
Must adjoining a predicate for any henselian valuation ring to an NIP field preserve NIP? The question uses the standard pure field and valuation-ring languages. The supplied residue-field-based literature assessment has a different hypothesis and requires clarification.
\subsection{Sources}
\begin{itemize}
\item Local UnsolvedMath registry: supplied record in \texttt{data/unsolvedmath/problems.json}, permanent ID \(30003023\), source problem code \texttt{catalogue entry}, title \emph{NIP Preservation under Henselian Valuations}. The record belongs to the Oberwolfach Reports Open Problems collection and gives proposed year 2016; no individual proposer is identified.
\item \emph{Model Theory: groups, geometry, and combinatorics} (2016), Oberwolfach Reports, pp.\ 5--50, DOI \texttt{10.4171/owr/2016/1}: \url{https://doi.org/10.4171/owr/2016/1}. This is the supplied report citation and also a supplied literature source. The background describes Oberwolfach Reports as documenting mathematical workshops and their open-problem discussions.
\item Supplied literature source labeled ``Anscombe, characterizing NIP henselian fields'': \url{https://doi.org/10.1112/jlms.12868}. No further bibliographic details are supplied.
\item Literature assessment checked at \texttt{2026-08-21}: the supplied curation labels the problem open and reports that the question concerning henselian expansions with NIP residue field in remaining imperfect or separably closed cases is unsettled, with broad special cases and structural classifications available. This is supplied dated curation derived from the OWR record, not an independently verified assessment; the record requests independent verification of time-sensitive status and cited claims.
\item The supplied statement status is \texttt{corrected\_verified}. Its verification note says that the statement was checked against the official Oberwolfach report and preserves that report's hypotheses, notation, and open target. This is a supplied source-verification claim, not verification performed for this draft.
\end{itemize}
\end{document}
